\section{Relaciones entre un anillo local noetheriano y su completado. Anillos excelentes}
\label{section:IV.7-fr}

Este apartado está dedicado principalmente a exponer ciertas
propiedades de los anillos noetherianos, generalmente estables
al pasar a un álgebra de tipo finito o al localizar, que se
verifican para los anillos que se encuentran con mayor
frecuencia (álgebras de tipo finito sobre $\mathbb Z$, sobre
un cuerpo o sobre un anillo local noetheriano completo), sin
que se verifiquen para todos los anillos noetherianos. Un
primer grupo de propiedades, relacionadas con la teoría de
la dimensión y, en particular, con la condición de las
cadenas, se estudia en los núms. 1 y 2; todas las propiedades
consideradas se verifican para los cocientes de anillos
noetherianos regulares, y sus demostraciones en este caso
son, en su mayoría, fáciles y bien conocidas [30]. Por ello,
para tales anillos, los desarrollos de los \S\S~1 y 2 carecen
de interés.

No ocurre lo mismo con las propiedades estudiadas en los
núms. 3, 4 y 5 (que no dependen de los núms. 1 y 2): sus
demostraciones, incluso para localizaciones de álgebras de
tipo finito sobre un cuerpo o sobre $\mathbb Z$ (donde se
deben en su mayoría a Zariski y Nagata), suelen ser
delicadas. Son ejemplos clásicos las equivalencias
$A$ normal $\Longleftrightarrow \widehat A$ normal y
$A$ reducido $\Longleftrightarrow \widehat A$ reducido.
Desarrollamos un método sistemático para formular y
demostrar tales propiedades mediante las propiedades de
las \emph{fibras} del morfismo canónico
$\operatorname{Spec}(\widehat A)\to\operatorname{Spec}(A)$
y utilizando las nociones y resultados expuestos en el \S~6.
El éxito de este método descansa, en último término, en
las excelentes propiedades de los \emph{anillos locales
completos} y de los anillos que se obtienen de ellos por
localización o al pasar a un álgebra de tipo finito. A este
respecto, desempeña un papel crucial el teorema de Nagata
(6.12.7), según el cual, para uno de estos anillos, el lugar
singular de $\operatorname{Spec}(A)$ es cerrado; lo mismo
ocurre, en lo relativo a las propiedades de permanencia,
con el criterio de regularidad (\textbf 0, 22.3.4), de naturaleza
más técnica.

Por último, en los núms. 6 y 7 aplicamos los resultados
obtenidos al estudio de la finitud de la clausura integral
de los anillos íntegros noetherianos.

El lector observará que los resultados más delicados del
\S~7 se utilizarán relativamente poco en el resto del
capítulo IV e incluso en los capítulos posteriores.
